\section*{Abstract}
\label{sec:abstract}

Automated chest X-ray (CXR) interpretation is vital for rapid clinical decision-making in emergency and inpatient settings. However, most existing automated baselines are image-centric, ignoring the clinical context—such as indications, vital signs, and prior imaging—that human clinicians routinely use. This thesis implements and evaluates a clinically-aligned, prior-aware multimodal CXR classification framework under practical computational constraints. The proposed model integrates current radiographs, free-text clinical indications, vital signs, and prior-study context using modality-specific encoders and a transformer-based fusion module. Evaluated on the long-tailed CXR-LT dataset, the proposed prior-aware multimodal model achieves a mean Average Precision (mAP) of 0.5595 and an AUROC of 0.9014, outperforming the image-only baseline (0.3246 mAP) and a standard multimodal baseline without prior information (0.4098 mAP). Notably, the integration of prior studies yields the largest improvement on rare tail classes, increasing mAP from 0.1586 to 0.3893. These findings demonstrate that incorporating prior-study context and multimodal clinical context significantly improves diagnostic performance and tail-class sensitivity under resource-constrained settings.
